\section{Neue CPU-Evaluation vom 14. November 2026}
\label{sec:neuevaluation}
Die neue Evaluation vergleicht drei gleich breite Additionskerne: den
Restoperator \texttt{\%}, nicht-negative Addition mit bedingter Subtraktion
und die kanonisch zentrierte $\pm B$-Korrektur. Alle Operanden sind 32 Bit
breit; $1\le B\le 2^{31}-1$ verhindert Überlauf in beiden Darstellungen.
Konstante Template-Moduli und Laufzeitmoduli werden getrennt ausgewertet.
Alle drei Methoden erhalten dieselben Compiler- und ISA-Optionen.

Der vollständige Lauf umfasst 630 Fälle und 12 Versuche
je Methode, mit allen sechs Ausführungsreihenfolgen gleich häufig und
einem vollständigen Warm-up. Fünf Eingabemuster decken positive, negative,
alternierende, grenznahe und pseudozufällige Schritte ab. Eine unabhängige
64-Bit-Referenz prüft alle Ausgaben und zählt die Korrekturen außerhalb
der Zeitmessung. Serielle Zustände, unabhängige Arrays und polynomielle
Addition werden getrennt betrachtet.

Gemessen wurde mit GCC 15.2.0, CMake Release, \texttt{-O3 -DNDEBUG}
und \texttt{-march=native} unter Windows auf dem Intel Core i9-14900K.
Der Prozess war an logische CPU 0 gebunden. Turbo, Temperatur und
Hintergrundlast wurden nicht kontrolliert. ARM und FPGA wurden nicht neu vermessen.
Die Zahlen gelten für diese Übersetzung und diesen Lauf.

\begin{center}\small
\begin{tabular}{llrrr}\toprule
Last & Bindung & Fälle & \% / zentriert & beschränkt / zentriert \\\midrule
serial & constant & 30 & 2.880 & 1.282 \\
serial & runtime & 30 & 5.144 & 0.962 \\
array & constant & 90 & 1.230 & 0.853 \\
array & runtime & 90 & 10.778 & 0.802 \\
polynomial & constant & 30 & 1.125 & 0.664 \\
polynomial & runtime & 30 & 11.791 & 0.677 \\
\bottomrule\end{tabular}\end{center}
Die Tabelle zeigt den Median der fallweisen Medianquotienten, ohne
die entartete Kontrolle $B=1$. Ein Quotient größer eins begünstigt die
zentrierte Darstellung. Diese deskriptive Zusammenfassung über unterschiedliche
Lasten ist kein gemeinsamer Anwendungsspeedup. Rohwerte und Quartile
jedes einzelnen Falls bleiben im Archiv erhalten; Quartile sind keine
Konfidenzintervalle. Verluste werden nicht entfernt.

Der Vergleich mit dem Restoperator allein überschätzt den Nachweis eines
Repräsentationsvorteils. Die beschränkte nicht-negative Baseline ist in
den zusammengefassten Array- und Polynomfällen schneller als die zentrierte
Variante. Bei seriellen konstanten Moduli zeigt die Zusammenfassung einen
Vorteil für centered; bei seriellen Laufzeitmoduli keinen allgemeinen Vorteil.
Die neue Messung stützt daher keine universelle Überlegenheit der Zentrierung.

Zusätzliche Fälle mit 65536 Zuständen und 1, 4, 16 oder 256 Updates
messen Normalisierung, Allokation und Rückkonvertierung mit. Eingabeerzeugung,
Validierung und abschließende Freigabe bleiben außerhalb. Diese Grenze ist
eine lokale Verarbeitungskette und kein vollständiges Kryptosystem.
Ein allgemeiner Amortisationspunkt wird nicht aus aggregierten Werten abgeleitet.

\noindent\textbf{Reproduktion:} Quellcode, Build-Kommandos, Hashes und alle
Rohdaten stehen im Repository unter \path{tests/results/x86/20261114_PAPER_V2}.
Die Methodik erläutert \path{docs/PAPER_BENCHMARK_V2_DE.md}.
Das Skript \path{scripts/run_paper_benchmarks.py} baut und misst die Matrix.
